\documentclass[11pt]{article}
\usepackage{amsmath,amssymb,amsthm,hyperref}
\newtheorem{theorem}{Theorem}
\newtheorem{lemma}[theorem]{Lemma}
\title{A Compact Cauchy Obstruction to Independent Pair Readout}
\author{Continuation checkpoint}
\date{23 November 2026}
\begin{document}
\maketitle

\paragraph{Scope.}
This note proves a limitation of one concrete permutation-memory mechanism.
It is a complete written proof, not a Lean-checked result and not a lower
bound for general attention heads. The original model and its denominator
admissibility conditions are unchanged. In particular, correlated subset
choices, signed residues, slot-dependent query poles, and more input choices
than output queries are not covered by the theorem below.

\begin{lemma}[A monotone bijection of a finite cube]
If a coordinatewise monotone map
$T:\{-1,1\}^k\longrightarrow\{-1,1\}^k$ is surjective, there is a
permutation $\sigma$ of $\{1,\ldots,k\}$ such that
$T_i(s)=s_{\sigma(i)}$ for every $i,s$.
\end{lemma}
\begin{proof}
Surjectivity between finite sets of equal size implies bijectivity.
As a permutation of a finite set, $T$ has finite order, say $T^m$ is the
identity. Its inverse $T^{m-1}$ is therefore also monotone. Thus $T$ is an
order automorphism. Identify the cube with the subsets of a $k$-element set,
ordered by inclusion. The empty set, the atoms (singleton sets), and least
upper bounds are defined by the order and hence preserved by an order
automorphism. The automorphism permutes the singleton sets. Every set is
the least upper bound of its singleton subsets, so its image is exactly the
union of their images. This is the claimed coordinate permutation.
\end{proof}

\begin{lemma}[Positive linear threshold bijections]
Let $V$ be a $k$ by $k$ matrix with strictly positive entries, and let
$\tau\in\mathbb R^k$. Suppose all $2^k$ strict sign patterns of
$Vs-\tau$ occur as $s$ ranges over $\{-1,1\}^k$. Then there is a
permutation $\sigma$ such that
\[
 V_{i,\sigma(i)}>\sum_{j\ne\sigma(i)}V_{ij}
 \quad\hbox{for every }i.
\]
\end{lemma}
\begin{proof}
There are only $2^k$ inputs. The assumed $2^k$ strict patterns therefore
use all inputs, each exactly once; no input has a zero output coordinate.
The sign map is coordinatewise monotone because every entry of $V$ is
positive. Apply the preceding lemma. For row $i$, put $j=\sigma(i)$.
The input with $s_j=1$ and every other coordinate equal to $-1$ gives
$V_{ij}-\sum_{\ell\ne j}V_{i\ell}>\tau_i$. The opposite input gives
$-V_{ij}+\sum_{\ell\ne j}V_{i\ell}<\tau_i$.
Combining the two strict inequalities proves the claim.
\end{proof}

\begin{theorem}[Compact independent-pair obstruction]
Fix $k\geq1$, $0\leq\beta<a\leq b$, and pairs
\[
 0\leq\lambda_j^-<\lambda_j^+\leq\beta
 \qquad(1\leq j\leq k).
\]
The pairs need not be mutually disjoint. For each input
$s\in\{-1,1\}^k$, select $\lambda_j^{s_j}$ independently, and query
\[
 L_i(s)=\sum_{j=1}^k\frac{1}{D_i-\lambda_j^{s_j}}-t_i,
 \qquad a\leq D_i\leq b,\quad 1\leq i\leq k,
\]
with arbitrary real thresholds $t_i$. If the $k$ queries realize every
strict sign pattern, then
\[
 k<1+\left(\frac{b}{a-\beta}\right)^2.
\]
In particular, when $a=1$, $b=11/10$, and $\beta=1/10$, this readout
cannot shatter $k\geq3$ independent pair choices.
\end{theorem}
\begin{proof}
Set $w_j=\lambda_j^+-\lambda_j^->0$. Subtract the midpoint of the two
possible values in each summand and absorb all midpoints into the
threshold. Multiplication by two does not change any sign, and gives a
threshold map with matrix
\[
 V_{ij}=\frac{1}{D_i-\lambda_j^+}
             -\frac{1}{D_i-\lambda_j^-}
 =\frac{w_j}{(D_i-\lambda_j^-)(D_i-\lambda_j^+)}.
\]
All denominators are strictly positive, and
\[
 \frac{w_j}{b^2}\leq V_{ij}\leq\frac{w_j}{(a-\beta)^2}.
\]
Write $W=\sum_jw_j$ and $\kappa=(b/(a-\beta))^2$.
By the preceding lemma there is a distinct dominant column
$j=\sigma(i)$ in every row. Its strict dominance and these bounds give
\[
 \frac{w_j}{(a-\beta)^2}>
 \frac{W-w_j}{b^2},\qquad
 w_j>\frac{W}{1+\kappa}.
\]
Every column occurs once. Summing gives
$W>kW/(1+\kappa)$; since $W>0$, this is $k<1+\kappa$.
For the specified compact interval,
$1+\kappa=1+121/81=202/81<3$.
\end{proof}

\paragraph{What the theorem rules out.}
One proposed hierarchical permutation code sends one member of each pole
pair into a slot block, and reads the choices by unweighted Cauchy sums at
nearby exterior poles. The theorem rules out reading one independent bit
per pair in such a square block once it has three or more pairs in the
specified compact regime. Large numbers of possible pair assignments do
not supply independently readable output bits.

\paragraph{Exact remaining gap.}
The theorem does not rule out a constant fraction of as many queries as
pair choices, or correlated equal-moment partitions used at overlapping
levels of a hierarchy. To obtain an asymptotic improvement from the latter
route one would need an explicit family of permutations of $H$ poles,
together with more than a constant multiple of $H$ fixed admissible
Cauchy queries, realizing every strict sign pattern. No such family is
constructed here. Mere factorial cardinality of the permutation family
does not establish this assertion.
\end{document}
